\section{Complete coefficient formulas and the first six degrees}\label{sec:coefficients}
The structural theorem becomes completely explicit by applying $D$ to the ordinary logarithm. Dynkin's formula is historically attributed to \cite{Dynkin}; modern computational treatments include \cite{CasasMurua}. The proof below derives its coefficients directly.

\subsection{Associative blocks and a finite formula for every word}
Set
\begin{equation}\label{eq:block-U}
 U=e^Xe^Y-1
   =\sum_{\substack{r,s\geq0\\r+s>0}}\frac{X^rY^s}{r!s!}.
\end{equation}
Substituting this into the logarithm gives the full associative expansion
\begin{equation}\label{eq:associative-all}
 Z=\sum_{k\geq1}\frac{(-1)^{k-1}}{k}
 \sum_{\substack{r_i,s_i\geq0\\r_i+s_i>0\ (1\leq i\leq k)}}
 \frac{X^{r_1}Y^{s_1}\cdots X^{r_k}Y^{s_k}}
 {\prod_{i=1}^k r_i!s_i!}.
\end{equation}
The inner sum is not a sum over independent commutators; at this stage it is simply noncommutative multiplication.

For a nonempty word $v$, define
\begin{equation}\label{eq:weight}
 a(v)=\begin{cases}
 1/(r!s!),&v=X^rY^s,\quad r+s>0,\\
 0,&\text{otherwise}.
 \end{cases}
\end{equation}
A cut $w=w_1\cdots w_k$ always means consecutive, nonempty subwords.

\begin{theorem}[All associative coefficients]\label{thm:word-cut}
For every word $w$ of length $n>0$,
\begin{equation}\label{eq:word-cut}
 c(w):=\coeff{w}Z
 =\sum_{k=1}^{n}\frac{(-1)^{k-1}}{k}
   \sum_{w=w_1\cdots w_k}\prod_{i=1}^k a(w_i).
\end{equation}
This is a finite, nonrecursive formula for every coefficient at every order. Equivalently,
\begin{equation}\label{eq:homogeneous-blocks}
 Z_n=\sum_{k=1}^{n}\frac{(-1)^{k-1}}{k}
 \sum_{\substack{r_i+s_i>0\\\sum_i(r_i+s_i)=n}}
 \frac{X^{r_1}Y^{s_1}\cdots X^{r_k}Y^{s_k}}
 {\prod_i r_i!s_i!}.
\end{equation}
\end{theorem}
\begin{proof}
The coefficient of $w$ in $U^k$ is the sum over all its cuts into $k$ nonempty pieces, with the product of their coefficients in $U$. Those coefficients are precisely $a(w_i)$. There are no contributions for $k>n$. Substitution in $\log(1+U)$ proves both formulas.
\end{proof}
For example, for $w=XXY$ the one-piece contribution is $1/2$, the two-piece contribution is $-(1+1/2)/2=-3/4$, and the three-piece contribution is $1/3$. Their sum is $1/12$. For $w=XYX$, the one-piece contribution is zero, the two-piece contribution is $-1/2$, and the three-piece contribution is $1/3$, giving $-1/6$.

\subsection{Dynkin's complete commutator formula}
\begin{theorem}[Dynkin formula]\label{thm:dynkin}
With $N=\sum_{i=1}^k(r_i+s_i)$,
\begin{equation}\label{eq:dynkin-all}
 \boxed{\displaystyle
 Z=\sum_{k\geq1}\frac{(-1)^{k-1}}{k}
 \sum_{\substack{r_i,s_i\geq0\\r_i+s_i>0}}
 \frac{R(X^{r_1}Y^{s_1}\cdots X^{r_k}Y^{s_k})}
 {N\prod_i r_i!s_i!}.}
\end{equation}
Equivalently,
\begin{equation}\label{eq:dynkin-word}
 Z_n=\frac1n\sum_{|w|=n}c(w)R(w).
\end{equation}
A summand in \cref{eq:dynkin-all} vanishes when $s_k>1$, or when $s_k=0$ and $r_k>1$.
\end{theorem}
\begin{proof}
Each $Z_n$ is a homogeneous Lie polynomial. Apply \cref{eq:dsw} to \cref{eq:homogeneous-blocks}; the operator replaces every word by its right-nested bracket and division by its length gives the factor $1/N$. Summing over degrees gives the formula. Under either stated condition the two innermost letters coincide, so the innermost bracket is zero.
\end{proof}
At every fixed degree the sums are finite. Analytic regrouping of all degrees requires a convergence argument; explicit sufficient conditions are proved in \cref{sec:convergence}.

\begin{corollary}[Every nonlinear term contains the basic commutator]\label{cor:comm-ideal}
For $n\geq2$, $Z_n$ belongs to the Lie ideal generated by $[X,Y]$. In particular it vanishes when $[X,Y]=0$.
\end{corollary}
\begin{proof}
In the quotient by that ideal, the two generators commute. Every bracket of length at least two in those generators is consequently zero, so the quotient Lie algebra is abelian. The image of $Z$ is $X+Y$ by the commuting exponential identity. Homogeneous components of higher degree therefore lie in the ideal.
\end{proof}
The phrase ``contains $[X,Y]$'' is an ideal-theoretic statement; it does not mean that an arbitrary chosen display must visibly have the same innermost bracket in every term.

\subsection{Symmetries, including a useful coefficient check}
\begin{proposition}\label{prop:symmetry}
The formal BCH series satisfies
\begin{align}
 Z(Y,X)&=-Z(-X,-Y),\label{eq:exchange}\\
 Z_n(Y,X)&=(-1)^{n+1}Z_n(X,Y),\label{eq:parity}\\
 \rev(Z_n(X,Y))&=(-1)^{n+1}Z_n(X,Y).\label{eq:reversal}
\end{align}
In particular, coefficients transform by the same sign under either swapping $X,Y$ or reversing every word.
\end{proposition}
\begin{proof}
Since $e^{-X}e^{-Y}=(e^Ye^X)^{-1}$, taking the formal logarithm gives \cref{eq:exchange}. Taking homogeneous components gives \cref{eq:parity}. Word reversal is an antihomomorphism, so it takes $e^Xe^Y$ to $e^Ye^X$ and commutes with the single-element logarithm series. This gives \cref{eq:reversal}.
\end{proof}

\subsection{Every term through degree six}
With the right-nesting convention \cref{eq:right-word}, the full answer through degree six is
\begin{align}
 Z_1={}&X+Y,\label{eq:Z1}\\
 Z_2={}&\tfrac12\Rword{XY},\label{eq:Z2}\\
 Z_3={}&\tfrac1{12}\bigl(\Rword{XXY}+\Rword{YYX}\bigr),\label{eq:Z3}\\
 Z_4={}&-\tfrac1{24}\Rword{YXXY},\label{eq:Z4}\\
 Z_5={}&-\tfrac1{720}\bigl(\Rword{YYYYX}+\Rword{XXXXY}\bigr)\nonumber\\
       &+\tfrac1{360}\bigl(\Rword{XYYYX}+\Rword{YXXXY}\bigr)\nonumber\\
       &+\tfrac1{120}\bigl(\Rword{YXYXY}+\Rword{XYXYX}\bigr),\label{eq:Z5}\\
 Z_6={}&\tfrac1{240}\Rword{XYXYXY}
       +\tfrac1{720}\bigl(\Rword{XYXXXY}-\Rword{XXYYXY}\bigr)\nonumber\\
       &+\tfrac1{1440}\bigl(\Rword{XYYYXY}-\Rword{XXYXXY}\bigr).\label{eq:Z6}
\end{align}
There is no omitted term of any of these degrees. Equations \cref{eq:word-cut,eq:dynkin-all} define every subsequent degree, rather than leaving the ellipsis unspecified.

For direct comparison with associative matrix multiplication,
\begin{align}
 Z_2={}&\tfrac12(XY-YX),\label{eq:assoc2}\\
 Z_3={}&\tfrac1{12}\bigl(X^2Y+XY^2-2XYX+Y^2X+YX^2-2YXY\bigr),\label{eq:assoc3}\\
 Z_4={}&\tfrac1{24}\bigl(X^2Y^2-2XYXY-Y^2X^2+2YXYX\bigr).\label{eq:assoc4}
\end{align}
For instance,
\[
 [X,[X,Y]]=X^2Y-2XYX+YX^2,
 \quad [Y,[Y,X]]=Y^2X-2YXY+XY^2,
\]
which proves the cubic conversion. Expanding $-[Y,[X,[X,Y]]]/24$ proves \cref{eq:assoc4}.

Here is a fully specified finite verification of the fifth- and sixth-degree displays, also usable for any larger display. For $w=a_1\cdots a_n$,
\begin{equation}\label{eq:right-expansion}
 R(w)=\sum_{S\subseteq\{1,\ldots,n-1\}}(-1)^{|S|}
 a_{i_1}\cdots a_{i_{n-1-|S|}}a_n
 a_{j_{|S|}}\cdots a_{j_1},
\end{equation}
where $i_1<\cdots<i_{n-1-|S|}$ enumerate $S^c$ and $j_1<\cdots<j_{|S|}$ enumerate $S$. At each bracket one either places the outer letter at the left or places it at the right with a minus sign, proving this identity by induction. Substitute \cref{eq:right-expansion} in \cref{eq:Z5,eq:Z6} and compare every word coefficient with the finite cut formula \cref{eq:word-cut}. The complete nonzero coefficient certificates are printed in \cref{sec:tables}; all remaining coefficients are zero in both calculations. This proves the displayed formulas by finite rational arithmetic, without choosing a commutator basis or appealing to a numerical approximation.
